\subsection{Objetivos particulares}
\label{Objetivos particulares}

\begin{itemize}
    \item \textbf{Implementar algoritmos de navegación:} Crear algoritmos avanzados que permitan al robot moverse de manera autónoma y eficiente en diferentes entornos, asegurando su capacidad para evitar obstáculos y alcanzar objetivos específicos.
    
    \item \textbf{Utilizar técnicas de IA para detectar objetivos:} Aplicar algoritmos de inteligencia artificial que permitan al robot identificar y localizar objetivos de manera precisa, optimizando su capacidad de respuesta en tiempo real.

    \item \textbf{Integrar sensores y actuadores en el prototipo:}
    Incorporar y coordinar diversos sensores y actuadores en el prototipo del robot, permitiéndole interactuar de manera precisa con su entorno y recolectar datos esenciales para su funcionamiento y control.
    
    \item \textbf{Realizar pruebas y validación en entornos simulados:}
    Desarrollar y ejecutar un programa de pruebas en el entorno virtual (Unreal Engine~5) para verificar, a alto nivel, el desempeño de la navegación autónoma (localización global, mapeo con LiDAR simulado y planificación de trayectorias), la respuesta del software de control ante eventos dinámicos de misión y la coherencia de la máquina de estados frente a escenarios cambiantes. Las métricas obtenidas en simulación guían la calibración de parámetros y la iteración de algoritmos antes de trasladarlos al prototipo físico.
    
    \item \textbf{Realizar pruebas y validación en el prototipo físico:}
    Llevar a cabo ensayos sobre el prototipo en condiciones controladas para evaluar, a alto nivel, la movilidad de la plataforma, el desempeño cinemático de la torreta, la capacidad de seguimiento de objetivos mediante visión computacional y la efectividad de neutralización, incluyendo escenarios con uno o más blancos en movimiento. Se verifica además la integración operativa de sensores, actuadores y subsistemas de navegación. Como objetivo de misión integrado, se busca validar la detección y eliminación de objetivos de aproximadamente $45 \times 45\ \text{cm}$ a una distancia de $20\ \text{m}$.

    
\end{itemize}